\begin{center}
\SourceLeaf{024b}\pdfbookmark[2]{6.}{article-006}6.
\end{center}
Progredimur iam ad comparationem motus cum \emph{tempore}. Statuendo ut in art. 1
spatium tempore $t$ circa Solem descriptum $= \frac{1}{2}g$, massam corporis moti $= \mu$,
posita massa Solis $= 1$, habemus $g = kt\BareRoot{p}.\BareRoot{(1+\mu)}$. Differentiale
spatii autem fit $= \frac{1}{2}rr\,\mathrm{d}v$, unde prodit $kt\BareRoot{p}.\BareRoot{(1+\mu)}
= \int rr\,\mathrm{d}v$, hoc integrali ita sumto, ut pro $t=0$ evanescat. Haec integratio
pro diversis sectionum conicarum generibus diverso modo tractari debet, quamobrem
singula iam seorsim considerabimus, initiumque ab ELLIPSI faciemus.

Quum $r$ ex $v$ per fractionem determinetur, cuius denominator e duabus partibus
constat, ante omnia hoc incommodum per introductionem quantitatis novae pro $v$
auferemus. Ad hunc finem statuemus $\operatorname{tang}\frac{1}{2}v\BareRoot{\frac{1-e}{1+e}}
= \operatorname{tang}\frac{1}{2}E$, quo pacto formula ultima art. praec. pro $r$ praebet
\[
r = \frac{p\cos\frac{1}{2}E^2}{(1+e)\cos\frac{1}{2}v^2} = p\left(\frac{\cos\frac{1}{2}E^2}{1+e}+\frac{\sin\frac{1}{2}E^2}{1-e}\right) = \frac{p}{1-ee}(1-e\cos E).
\]
Porro fit $\frac{\mathrm{d}E}{\cos\frac{1}{2}E^2} = \frac{\mathrm{d}v}{\cos\frac{1}{2}v^2}\BareRoot{\frac{1-e}{1+e}}$,
adeoque $\mathrm{d}v = \frac{p\,\mathrm{d}E}{r\BareRoot{(1-ee)}}$; hinc
\[
rr\,\mathrm{d}v = \frac{r\,\grave{p}\,\mathrm{d}E}{\BareRoot{(1-ee)}} = \frac{pp}{(1-ee)^{\frac{3}{2}}}(1-e\cos E)\,\mathrm{d}E,
\]
atque integrando
\[
kt\BareRoot{p}.\BareRoot{(1+\mu)} = \frac{pp}{(1-ee)^{\frac{3}{2}}}(E-e\sin E)+\mathrm{Const.}
\]
Quodsi itaque tempus a transitu per perihelium inchoamus, ubi $v=0$, $E=0$
\SourceLeaf{025}adeoque Const. $= 0$, habebimus, propter $\frac{p}{1-ee}=a$,
\[
E-e\sin E = \frac{kt\BareRoot{(1+\mu)}}{a^{\frac{3}{2}}}\cdot
\]

In hac aequatione angulus auxiliaris $E$, qui \emph{anomalia excentrica} dicitur, in
partibus radii exprimi debet. Manifesto autem hunc angulum in gradibus etc. retinere
licet, si modo etiam $e\sin E$ atque $\frac{kt\BareRoot{(1+\mu)}}{a^{\frac{3}{2}}}$ eodem modo
exprimantur; in minutis secundis hae quantitates exprimentur, si per numerum
$206264{,}806$ multiplicantur. Multiplicatione quantitatis posterioris supersedere
possumus, si statim quantitatem $k$ in secundis expressam adhibemus, adeoque, loco
valoris supra dati, statuimus $k = 3548^{\prime\prime}{,}18761$, cuius logarithmus $= 3{,}550\,0065\,746$. —
Hoc modo expressa quantitas $\frac{kt\BareRoot{(1+\mu)}}{a^{\frac{3}{2}}}$ \emph{anomalia
media} vocatur, quae igitur in ratione temporis crescit, et quidem quotidie augmento
$\frac{k\BareRoot{(1+\mu)}}{a^{\frac{3}{2}}}$, quod \emph{motus medius diurnus} dicitur.
Anomaliam mediam per $M$ denotabimus.
